\section{Introduction}
\label{sec:introduction}

Quantum algorithms for scientific computing are commonly specified through
access models: a block-encoding of a matrix, sparse location and entry
oracles, or a state preparation for an input vector.
This viewpoint connects quantum linear-system solvers
\cite{harrow2009hhl,costa2022discrete}, Hamiltonian simulation
\cite{berry2015taylor,childs2021trotter,low2019qubitization},
quantum singular-value transformation \cite{gilyen2019qsvt,martyn2021grand},
and differential-equation algorithms
\cite{an2026lchs,jin2024schrodingerization,liu2021pnas}.
Implementing such a paper requires decisions that its oracle-level
description deliberately leaves open: how data are accessed, how matrix
entries are computed reversibly, which finite-precision representation is
used, and how intermediate results and success signals are consumed.

Our primary user is the quantum algorithm researcher implementing,
modifying, and comparing these constructions. This user needs to preserve
the algorithm's subroutine structure while progressively replacing its
access assumptions with executable implementations. A useful framework
must support both the abstract program and its concrete realizations,
connect numerical checks to the implementation being checked, and expose
the resource consequences of implementation choices.

We present \pqlang{}, a Python framework for implementing and analyzing
quantum scientific-computing algorithms. Its embedded programming interface
generates a modular register-level intermediate representation (RIR).
An algorithm can be stored with declared but unimplemented oracle modules,
restored independently of its Python generator, and partially linked to
concrete operations. Register signatures, views, calls, static repetition,
control, and adjoint remain explicit until backend lowering. RIR widths
and assembly constants are concrete; deferred binding concerns module
implementations, not arbitrary symbolic parameter substitution.

The framework connects four parts of the research workflow:
\begin{enumerate}
\item \textbf{Express access requirements and compose algorithms.}
  Typed role views describe block-encodings, state preparations, and sparse
  access. Python structural interfaces let application objects supply these
  views without inheriting framework base classes. Algorithm contracts
  check the resolved views, and the generator uses the same objects that
  were checked. Structural compatibility and mathematical assumptions
  remain distinct.
\item \textbf{Choose realizations through partial linking.}
  Open modules are serialized program objects. Binding checks interface,
  scale, and capability compatibility, preserves unresolved dependencies,
  and connects captured resources through the module graph. Binding
  reports record program fingerprints, resource mappings, and remaining
  obligations. A change to an assembly-time constant such as a
  block-encoding scale requires regeneration.
\item \textbf{Construct scientific-computing oracles.}
  QRAM access and reversible arithmetic provide complementary
  implementations of access models. A compiler for supported pure Python
  numerical functions lowers mathematical formulas through a typed
  intermediate representation to reversible fixed-point modules with
  explicit rounding and status conventions. Applications can load
  precomputed entries or compute entries from raw data.
\item \textbf{Verify behavior and compare costs.}
  Small-instance simulations are compared with independent mathematical
  or finite-precision references. Open-program analysis reports known
  circuit costs together with unresolved oracle calls; closed-program
  analysis follows the implemented lowering recipes. QRAM accesses are
  recorded by logical memory resource, and repeated rotations are counted
  without expanding their occurrences.
\end{enumerate}

The contribution is this implemented workflow and its supporting
interfaces, linking mechanisms, and scientific-computing components.
We do not claim novelty for structural typing, modular intermediate
representations, or resource counting individually. Qualtran
\cite{harrigan2024qualtran} and pyLIQTR \cite{obenland2026pyliqtr}
are close precedents for compositional fault-tolerant algorithm
implementation and analysis. Our comparisons therefore examine concrete
research tasks, access-model realizations, and artifact boundaries, rather
than treating a new programming syntax or a longer algorithm list as
sufficient evidence of progress.

The implementation includes block-encoding algebra, signal-processing
transforms, state preparation, and linear-system and differential-equation
constructions. A controlled oracle-realization study isolates the effect
of gate-table, QRAM-table, and arithmetic implementations on one stored
algorithm. Roe/QFVM and QHAM case studies exercise deeper scientific
compositions. Verification results distinguish structural checks,
finite-instance numerical agreement, backend agreement, and unresolved
accuracy or convergence assumptions. These are complementary evidence,
not a proof of correctness for every parameter choice.

\paragraph{Organization.}
Section~\ref{sec:requirements} develops the requirements of this research
workflow, and Section~\ref{sec:related} places the framework among existing
systems. Sections~\ref{sec:language} and \ref{sec:oracles} describe RIR,
interfaces, resources, and linking. Sections~\ref{sec:primitives} and
\ref{sec:ode} present the reusable computations and solver constructions.
Section~\ref{sec:casestudies} compares implementations and presents
scientific applications. Section~\ref{sec:implementation} reports
verification and resource analysis, with detailed inventories and cost
studies in the appendices. Section~\ref{sec:discussion} states limitations
and subsequent research directions.
